\section{Incidence correspondance and the universal Fano scheme}

We aim to introduce the powerful concept of incidence correspondences, which will be a tool we will use several times later in this thesis. 

\begin{definition}
    (\cite{EH16} Section 6.1). N = $\binom{n+d}{d}-1$

    \[\Phi = \{(X, L) \in \mathbb{P}^N \times \mathbb{G}(k,n) \mid L \subseteq X\}\]
\end{definition}

\begin{proposition}
    $\Phi \subset \mathbb{P}^N \times \mathbb{G}(k,n)$ is a closed subset of $\mathbb{P}^N \times \mathbb{G}(k,n)$.
\end{proposition}

\todo[inline]{Decide if we want to prove, or refer to proof. Can refer to Harris 1995 chapter 6, or perhaps proof of 22.31 in EO}

\begin{proposition}[EH16, Proposition 6.1]
    Let $N = \binom{n+d}{d}-1$. The universal Fano scheme $\Phi = \Phi(n,d,k) \subseteq \mathbb{P}^N \times \mathbb{G}(k,n)$ is a smooth irreducible variety of dimension
    \[\dim \Phi(n,d,k) = N 0 (k+1)(n-k) - \binom{k+d}{k}.\]
\end{proposition}

\begin{corollary}[EH16, corollary 6.2]
    \todo[inline]{Include result.}
\end{corollary}

\begin{corollary} \label{cor:Fano}
    (\cite{EO25}, 22.31) For a hypersurface $S_d \subseteq \mathbb{P}^3$ of degree $d$, we expect to find no lines on $S_d$ when $d \geq 4$, finitely many lines when $d=3$, and positive-dimensional families of lines when $d \leq 2$.
\end{corollary}